
\section{Case study: one candidate end to end}\label{app:case}
This appendix traces the top screen hit through every stage, collecting
numbers reported piecemeal above into one falsifiable narrative.

\paragraph{Stage 1: pool.} The discovery pool is 852 reviewed,
uncharacterized UniProt proteins of 10--60 residues, fetched with
provenance logging. P0ACW4 (YdcA, \emph{E. coli} K-12, 57 aa) enters as an
uncharacterized protein; the screen makes no use of any AMP annotation.

\paragraph{Stage 2: scoring.} The frozen study-07 ensemble scores it
0.9082 (CNN 0.9992, GNN 0.8171), the top unique score in the pool. The
duplicate-accession pair P0ACW4/P0ACW5 carries an identical 57-residue
sequence in two species records and is reported as such.

\paragraph{Stage 3: novelty audit.} Maximum 3-mer Jaccard against APD6
2024 is 0.062 (novelty 0.938) and against the KW-0929 corpus 0.055; the
longest common substring with any APD6 entry is 6 residues. Under
Definition of Eq.~\eqref{eq:novelty} this is weak local similarity only.

\paragraph{Stage 4: physicochemical audit.} Length 57, net charge
$+4.87$ at pH 7.4 (Eq.~\eqref{eq:netcharge}), mean hydropathy $+0.08$,
Eisenberg moment $1.18$ (Eq.~\eqref{eq:hmoment}), Boman index $1.09$:
cationic and roughly hydrophobicity-neutral, a plausible but not extreme
AMP-like profile.

\paragraph{Stage 5: segment ablation.} YdcA carries an annotated 20-aa
N-terminal signal peptide (EcoCyc/EcoliWiki gene record). Scoring the
segments separately with the same frozen ensemble: full precursor 0.9082;
precursor minus signal segment 0.6299. The model's confidence is carried
substantially by the signal segment, which is expected: signal peptides
are hydrophobic, and the confounded training set rewards hydrophobic
cationic stretches (Section~\ref{sec:labelconfound}).


\begin{table}[h]\centering\small
\caption{Signal-peptide ablation for P0ACW4 with the frozen ensemble
(study 15; annotation: EcoliWiki ydcA gene record, signal residues 1--20).}
\label{tab:signalabl}
\begin{tabular}{lcccc}
\hline
Segment & Length & CNN & GNN & Ensemble \\
\hline
Full precursor (1--57) & 57 & 0.9992 & 0.8171 & 0.9082 \\
Signal segment (1--20) & 20 & 0.5905 & 0.7920 & 0.6912 \\
Post-signal segment (21--57) & 37 & 0.9974 & 0.2623 & 0.6299 \\
\hline
\end{tabular}
\end{table}

\paragraph{Verdict.} P0ACW4 is a named, quantified, falsifiable
\emph{hypothesis}: broth-microdilution MIC against \emph{E. coli} K-12 and
\emph{S. aureus} ATCC 25923 decides it. The ablation is precisely the kind
of negative control that keeps this report honest: the high score partly
reflects a trafficking signal, not necessarily antimicrobial chemistry.
